\section{Autodualidad de la unidad interna de Planck bajo inversión de orientación}

El selector CT108 construye

\begin{equation}
\theta_{108}=\left(\frac{54}{\pi}\right)^2,
\qquad
\ell_P=\frac{54}{\pi}\ell_0,
\qquad
t_P=\frac{54}{\pi}t_0,
\label{eq:planck-length-time}
\end{equation}

y

\begin{equation}
E_P=\frac{\hbar_{\rm int}}{t_P}
=\frac{\pi\hbar_{\rm int}}{54t_0}.
\label{eq:planck-energy}
\end{equation}

Comparando \cref{eq:boltzmann-clock,eq:planck-energy} se obtiene

\begin{equation}
\boxed{E_P=k_B\Theta_{\rm clk}\log3.}
\label{eq:planck-trit}
\end{equation}

La relaci\'on no depende de una coincidencia entre unidades: ambos miembros
proceden de la misma estructura cronol\'ogica \(108\).

Una vez derivada internamente la constante gravitatoria en el
\cref{ch:gravity}, la familia de Planck satisface

\begin{align}
\ell_P^2&=\frac{G\hbar}{c^3}=\theta_{108}\ell_0^2,
\label{eq:planck-area}\\
F_P&=\frac{c^4}{G},
\qquad
P_P=\frac{c^5}{G},
\label{eq:planck-force-power}\\
\rho_P&=\frac{\hbar}{\theta_{108}^2c^5t_0^4}.
\label{eq:planck-density}
\end{align}

Aunque \(G\) y \(\hbar\) se transformen rec\'iprocamente entre las secciones
anterior y posterior al retorno, su producto en \(\ell_P^2\) conserva la
geometr\'ia. Esta invariancia determina el \`area como interfaz entre las
realizaciones geom\'etrica, informacional y de entrelazamiento.
